% !TEX root = ../slides.tex
%%%%%%%%%%%%%%%%%%%%%%%%%%%%%%%%%%%%%%%%%%%%%%%%%%%%%%%%%%%%%%%%%%%%%%%%%%%%%%%%
% Session 5: Finite-elements spaces
%%%%%%%%%%%%%%%%%%%%%%%%%%%%%%%%%%%%%%%%%%%%%%%%%%%%%%%%%%%%%%%%%%%%%%%%%%%%%%%%

\session{session05}{The finite-elements method}{Finite-elements spaces}
{Sections 2.3--2.4}
{--}
{2.7--2.10}

\begin{frame}{Plan for today}
    Road map of the finite-element method:
    \begin{enumerate}
        \item \textbf{One piece}: $(K, \mathcal P, \mathcal N)$, nodal basis, interpolant $\mathcal I_K$. \hfill Session 4
        \item \textbf{The error on one piece}: $[v - \mathcal I_K v]_{H^1(K)} \le \frac{C}{\mu_K} h_K [v]_{H^2(K)}$. \hfill Session 4
        \item \textbf{Many pieces}: glue them to build $V_h$. \hfill {\color{mnuaccent}today}
        \item \textbf{The global error}: add the errors of all the pieces. \hfill {\color{mnuaccent}today}
    \end{enumerate}
    \bigskip
    Today:
    \begin{itemize}
        \item sub-divisions and triangulations; the finite-element space;
        \item glueing: when is a piece-wise function in $H^1$?
        \item the interpolation error in the whole domain;
        \item the convergence of the finite-element method with linear elements.
    \end{itemize}
\end{frame}

\begin{frame}{Sub-divisions and triangulations}
    \begin{definition}[Sub-division]
        A sub-division of $\Omega$ is a finite collection $\mathcal T$ of open sets $K$ such that
        \[
            \bigcup_{K \in \mathcal T} \overline K = \overline \Omega, \qquad K \cap K' = \emptyset \text{ if } K \ne K'.
        \]
    \end{definition}
    \begin{definition}[Triangulation]
        Let $\Omega \subset \R^2$ be polygonal. A triangulation is a sub-division into triangles such that
        no vertex of a triangle lies in the interior of an edge of another triangle.
    \end{definition}
    \begin{center}
        \femfigure[.8]{triangulation}
    \end{center}
    Equivalently: $\overline K \cap \overline{K'}$ is empty, a common vertex, or a complete common edge.
\end{frame}

\begin{frame}{Finite-element space}
    \small
    \begin{definition}[Finite-element space]
        Given a sub-division $\mathcal T$ and a finite element $(K, \mathcal P_K, \mathcal N_K)$ on each $K \in \mathcal T$,
        \[
            \mathcal P_{\mathcal T} \defeq \prod_{K \in \mathcal T} \mathcal P_K ,
        \]
        identified with piece-wise functions defined almost everywhere in $\Omega$.
    \end{definition}
    \begin{definition}[Global interpolant and global nodal functions]
        \[
            (\mathcal I_{\mathcal T} v)\big|_K = \mathcal I_K (v|_K),
            \qquad
            N_{K,i}(v) = N_i^{(K)}(v|_K),
            \qquad K \in \mathcal T
        \]
        (values on $\partial K$ are taken as limits from inside $K$). We often write $\mathcal T_h$ and $\mathcal I_h$.
    \end{definition}
    \begin{itemize}
        \item \alert{No continuity} is assumed between elements: $\mathcal I_{\mathcal T} v$ may jump across their boundaries.
        \item But Galerkin's method needs $V_h \subset V$ (e.g.\ $V = H^1$): this requires some continuity.
    \end{itemize}
\end{frame}

\begin{frame}{Example in $d=1$: two intervals}
    $\Omega = (-1,1)$, \ $\mathcal T = \{(-1,0), (0,1)\}$, linear elements.
    \medskip

    \begin{columns}[T]
        \begin{column}{.5\textwidth}
            Four local basis functions:
            \[
                \varphi_1 = -x, \ \varphi_2 = 1 + x \text{ on } (-1,0),
            \]
            \[
                \varphi_3 = 1 - x, \ \varphi_4 = x \text{ on } (0,1)
            \]
            (extended by $0$). $\dim \mathcal P_{\mathcal T} = 4$.
            \medskip

            Nodal functions: $v(-1^+)$, $v(0^-)$, $v(0^+)$, $v(1^-)$.
        \end{column}
        \begin{column}{.45\textwidth}
            \centering
            \begin{tikzpicture}[scale=1.3]
                \draw[-{Stealth}] (-1.2,0) -- (1.3,0);
                \draw[thick,blue] (-1,1) -- (0,0);
                \draw[thick,blue!50!cyan] (-1,0) -- (0,1);
                \draw[thick,red!70!black] (0,1) -- (1,0);
                \draw[thick,orange] (0,0) -- (1,1);
                \node[blue,left] at (-1,1) {\scriptsize $\varphi_1$};
                \node[blue!50!cyan] at (-.3,1.15) {\scriptsize $\varphi_2$};
                \node[red!70!black] at (.3,1.15) {\scriptsize $\varphi_3$};
                \node[orange,right] at (1,1) {\scriptsize $\varphi_4$};
                \foreach \x in {-1,0,1} \fill (\x,0) circle (.8pt) node[below] {\scriptsize $\x$};
            \end{tikzpicture}
        \end{column}
    \end{columns}
    \medskip
    Continuous functions:
    \[
        \mathcal P_{\mathcal T} \cap H^1(-1,1) = \mathcal P_{\mathcal T} \cap C([-1,1]) = \spannedby\{\varphi_1,\ \varphi_2 + \varphi_3,\ \varphi_4\}.
    \]
    Only three nodal functions survive: $v(-1)$, $v(0)$, $v(1)$. The hat functions are the dual basis.
\end{frame}

\begin{frame}{Example in $d=2$: two triangles}
    \begin{columns}[c]
        \begin{column}{.62\textwidth}
            $\Omega = (0,1)^2$, \ $\mathcal T = \{K_1, K_2\}$ sharing the edge $E$, linear triangles.
            \begin{itemize}
                \item $\mathcal P_{\mathcal T}$: a polynomial of degree $\le 1$ in each triangle. \ $\dim \mathcal P_{\mathcal T} = 3 + 3 = 6$.
                \item Six nodal functions: the vertices $(1,0)$ and $(0,1)$ count twice.
            \end{itemize}
        \end{column}
        \begin{column}{.34\textwidth}
            \centering
            \femfigure[.75]{two-triangles}
        \end{column}
    \end{columns}
    \medskip
    \textbf{When is $v \in \mathcal P_{\mathcal T}$ continuous?} Along $E$, both pieces are linear functions of one variable: they coincide on $E$ iff they coincide at its two end points.
    \[
        \dim\big(\mathcal P_{\mathcal T} \cap C(\overline\Omega)\big) = 6 - 2 = 4 = \text{number of vertices}.
    \]
    Dual basis: one continuous function per vertex, glueing the local basis functions. For example
    \[
        \varphi_{(0,0)} = \begin{cases} 1 - x - y & \text{in } K_1, \\ 0 & \text{in } K_2, \end{cases}
        \qquad
        \varphi_{(1,0)} = \begin{cases} x & \text{in } K_1, \\ 1 - y & \text{in } K_2. \end{cases}
    \]
\end{frame}

\begin{frame}{Glueing theorem}
    \small
    \begin{theorem}[Glueing]
        Let $\overline\Omega = \overline{\omega_1} \cup \overline{\omega_2}$, $\omega_1 \cap \omega_2 = \emptyset$ (Lipschitz boundaries), $u_i \in C^1(\overline{\omega_i})$, and
        \[
            u = u_i \text{ in } \omega_i, \qquad g_k = \tfrac{\partial u_i}{\partial x_k} \text{ in } \omega_i.
        \]
        Then, for $\varphi \in C^\infty_c(\Omega)$,
        \[
            \int_\Omega u \frac{\partial \varphi}{\partial x_k} = -\int_\Omega g_k \varphi + \int_{\partial\omega_1 \cap \partial\omega_2} (u_1 - u_2)\, \varphi\, n_{\omega_1}\cdot e_k .
        \]
        In particular, $u \in H^1(\Omega)$ if and only if $u_1 = u_2$ on $\partial\omega_1 \cap \partial\omega_2$.
    \end{theorem}
    \begin{proofsteps}[Idea]
        Integrate by parts in $\omega_1$ and in $\omega_2$ separately: the boundary terms on the common interface do not cancel unless the traces agree.
    \end{proofsteps}
    \takeaway{Piece-wise smooth functions are in $H^1$ iff they are continuous.}
\end{frame}

\begin{frame}{Choosing $V_h$}
    If the solution $u \in C^r(\overline \Omega)$, the natural space is
    \[
        V_h \defeq V \cap \mathcal I_h C^r(\overline\Omega) \ \subset\ V \cap \mathcal P_{\mathcal T}
    \]
    (usually with equality). Then $\mathcal I_h u \in V_h$ and we can use Céa's lemma.
    \begin{example}[Linear elements]
        \begin{itemize}
            \item Dirichlet Laplacian: $V_h = H^1_0(\Omega) \cap \mathcal P_{\mathcal T}$: continuous, piece-wise linear, zero on $\partial\Omega$.
            \item Helmholtz with Neumann conditions: $V_h = H^1(\Omega) \cap \mathcal P_{\mathcal T}$.
        \end{itemize}
    \end{example}
    \textbf{Building a basis of $V_h$.}
    \begin{itemize}
        \item Glue the local basis functions that share a degree of freedom ($N_{K,i} = N_{L,j}$ on neighbouring elements); impose the boundary conditions.
        \item In practice: a \alert{local-to-global} map, and matrices computed element by element (Sessions 6--7).
        \item If every element is affine equivalent to $(\widehat K, \widehat{\mathcal P}, \widehat{\mathcal N})$, we call it the \emph{reference element}.
    \end{itemize}
\end{frame}

\begin{frame}{The global error: adding the local errors}
    How well does $\mathcal I_{\mathcal T} v$ approximate $v$ in the whole domain?
    \medskip

    Integrals over $\Omega$ are sums of integrals over the elements: for $w \in H^1(\Omega)$,
    \[
        \|w\|^2_{L^2(\Omega)} = \sum_{K \in \mathcal T} \|w\|^2_{L^2(K)},
        \qquad
        [w]^2_{H^1(\Omega)} = \sum_{K \in \mathcal T} [w]^2_{H^1(K)} .
    \]
    Apply it to $w = v - \mathcal I_{\mathcal T} v$, which in each $K$ is the local error $v - \mathcal I_K v$ (Session 4):
    \[
        [v - \mathcal I_K v]_{H^1(K)} \le \frac{C}{\mu_K}\, h_K\, [v]_{H^2(K)} .
    \]
    Notation for a sub-division $\mathcal T$:
    \[
        h = \max_{K \in \mathcal T} h_K \ \ (\text{the \emph{mesh size}}),
        \qquad
        \mu_* = \min_{K \in \mathcal T} \mu_K = \min_{K \in \mathcal T} \frac{\rho_K}{h_K}.
    \]
\end{frame}

\begin{frame}{Interpolation error for linear elements}
    \begin{theorem}
        Let $\Omega$ be an interval ($d=1$) or a polygon ($d=2$), $\mathcal T$ a sub-division into intervals or a triangulation, with the linear Lagrange element in each $K$.
        Then, for every $v \in H^2(\Omega)$, $\mathcal I_{\mathcal T} v$ is continuous, belongs to $H^1(\Omega)$, and
        \[
            \|v - \mathcal I_{\mathcal T} v\|_{L^2(\Omega)} \le C h^2\, [v]_{H^2(\Omega)},
            \qquad
            [v - \mathcal I_{\mathcal T} v]_{H^1(\Omega)} \le \frac{C}{\mu_*}\, h\, [v]_{H^2(\Omega)} .
        \]
    \end{theorem}
    \begin{proofsteps}
        \begin{enumerate}
            \item $v$ is continuous ($d \le 2$); neighbouring interpolants share the values at the common vertices: $\mathcal I_{\mathcal T} v$ is continuous, hence in $H^1(\Omega)$ (glueing).
            \item Add the local errors:
                  \[
                      [v - \mathcal I_{\mathcal T} v]^2_{H^1(\Omega)} = \sum_{K} [v - \mathcal I_K v]^2_{H^1(K)}
                      \le \sum_K \frac{C^2}{\mu_K^2} h_K^2 [v]^2_{H^2(K)}
                      \le \frac{C^2}{\mu_*^2} h^2 [v]^2_{H^2(\Omega)} .
                  \]
        \end{enumerate}
    \end{proofsteps}
\end{frame}

\begin{frame}{Convergence of the finite-element method}
    \small
    \begin{theorem}[Linear elements for the Dirichlet Laplacian]
        Let $\Omega \subset \R^2$ be a polygon, $f \in L^2(\Omega)$, and $u \in H^1_0(\Omega)$ the weak solution of $-\Delta u = f$. Assume $u \in H^2(\Omega)$.
        Let $\mathcal T$ be a triangulation, $V_h = \mathcal P_{\mathcal T} \cap H^1_0(\Omega)$ and $u_h \in V_h$ the Galerkin solution. Then
        \[
            \|u - u_h\|_{H^1(\Omega)} \le \frac{C}{\mu_*}\, h\, [u]_{H^2(\Omega)} .
        \]
    \end{theorem}
    \begin{proofsteps}
        \begin{enumerate}
            \item Céa: $\|u - u_h\|_{H^1} \le \frac{M}{\alpha^*} \|u - v_h\|_{H^1}$ for all $v_h \in V_h$.
            \item Take $v_h = \mathcal I_{\mathcal T} u \in V_h$: continuous, piece-wise linear, and zero on $\partial\Omega$ (zero at the boundary vertices, linear on each edge).
            \item Interpolation error: $\|u - \mathcal I_{\mathcal T} u\|_{H^1} \le C\big(h^2 + \frac{h}{\mu_*}\big) [u]_{H^2}$.
        \end{enumerate}
    \end{proofsteps}
    \begin{itemize}
        \item $u \in H^2(\Omega)$ holds if $\Omega$ is convex. Then Aubin--Nitsche gives $\|u - u_h\|_{L^2} \le \frac{C}{\mu_*^2} h^2 \|f\|_{L^2}$.
    \end{itemize}
\end{frame}

\begin{frame}{Shape regularity and mesh generation}
    \begin{columns}[c]
        \begin{column}{.48\textwidth}
            The constants blow up if the triangles become flat ($\mu_K \to 0$, Exercise 2.9).
            \begin{definition}[Shape-regular family]
                A family of triangulations with
                \[
                    \mu_* \ge \mu > 0, \qquad \mu \text{ independent of } h.
                \]
            \end{definition}
        \end{column}
        \begin{column}{.5\textwidth}
            \centering
            \femfigure[.55]{h-rho}
        \end{column}
    \end{columns}
    \takeaway{For shape-regular triangulations: error $O(h)$ in $H^1$ and $O(h^2)$ in $L^2$.}
    \bigskip
    \textbf{Mesh generation.} Given points in the plane, how do we choose the triangles?
    \begin{itemize}
        \item We must control the shape ratio $\mu_K = \rho_K/h_K$.
        \item A popular choice: Delaunay triangulations (e.g.\ the program \texttt{Triangle}).
    \end{itemize}
\end{frame}

\begin{frame}{Other elements and other norms}
    The same argument (add the local errors) works in general.
    \begin{theorem}[Affine-equivalent sub-divisions]
        Assume that $\mathcal I_{\mathcal T}: H^m(\Omega) \to H^k(\Omega) \cap \mathcal P_{\mathcal T}$, every element is affine equivalent to the same reference element $(\widehat K, \widehat{\mathcal P}, \widehat{\mathcal N})$, and
        $[\widehat v - \mathcal I_{\widehat K} \widehat v]_{H^k(\widehat K)} \le \widehat C\, [\widehat v]_{H^m(\widehat K)}$. Then
        \[
            [v - \mathcal I_{\mathcal T} v]_{H^k(\Omega)} \le C \sup_{K \in \mathcal T} \frac{h_K^{m-k}}{\mu_K^k}\, [v]_{H^m(\Omega)} \qquad \forall v \in H^m(\Omega).
        \]
    \end{theorem}
    \begin{itemize}
        \item Elements containing the polynomials of degree $m-1$: interpolation error $\lesssim h^{m-k}$ in $H^k$, and by Céa the Galerkin error is $\lesssim h^{m-1}$ in $H^1$.
        \item Without affine equivalence (e.g.\ Hermite elements): the same rate, using the Bramble--Hilbert estimate on each element, for shape-regular meshes.
    \end{itemize}
\end{frame}

\begin{frame}{Summary and practice}
    \begin{itemize}
        \item $\mathcal P_{\mathcal T} = \prod_K \mathcal P_K$; global interpolant defined element by element.
        \item Glueing: piece-wise smooth $+$ continuous $\Rightarrow H^1$. Hence $V_h = V \cap \mathcal P_{\mathcal T}$.
        \item Global error $=$ sum of local errors: $[v - \mathcal I_{\mathcal T} v]_{H^1(\Omega)} \le \frac{C}{\mu_*} h [v]_{H^2(\Omega)}$ for linear elements.
        \item With Céa: the finite-element method converges with order $h$ in $H^1$, on shape-regular meshes.
    \end{itemize}
    \bigskip
    \textbf{Practice.} Exercises 2.7--2.10:
    \begin{itemize}
        \item 2.7: a sub-division with a hanging vertex, and how to refine it.
        \item 2.8: hat functions on a triangulation form a basis; effect of Dirichlet conditions.
        \item 2.9: degenerate triangles and the shape ratio.
        \item 2.10: a continuous piece-wise linear function with a kink: in $H^1$? in $H^2$?
    \end{itemize}
    \bigskip
    \textbf{Next session:} finite elements in $d = 1$.
\end{frame}
